% Generado por veredicto.py a partir de plan2/. No editar a mano.
\begin{table}[H]
\setstretch{1}
\centering
\scriptsize
\setlength{\tabcolsep}{1.5pt}
\begin{tabular}{llcrrrrrrcl}
\toprule
Caso & \makecell[l]{Red y\\demanda} & Sem. & Carr. & \makecell{Dem.\\(veh/h)} & \makecell{Retraso\\(s/veh)} & \makecell{Cambio\\frente al\\fijo} & \makecell{Cambio\\frente al\\actuado} & \makecell{Cambio de\\la recom-\\pensa} & Conv. & Veredicto \\
\midrule
P1 & \makecell[l]{Corredor,\\{}nominal} & 2 & 1 & 1810 & \makecell[r]{23.64\\{}$\pm$ 0.15} & \makecell[r]{+9.1\,\%\\{}[7.7,\,10.3]\\$p$ $<$\,0.001} & \makecell[r]{+10.2\,\%\\{}[9.3,\,11.1]\\$p$ $<$\,0.001} & \makecell[r]{$-$11.6\,\%\\{}[$-$13.0,\,$-$10.0]\\$p$ $<$\,0.001} & \makecell[c]{sí\\{}(1/3)} & FALLA \\
P2 & \makecell[l]{Corredor,\\{}alta} & 2 & 1 & 3620 & \makecell[r]{498.53\\{}$\pm$ 11.50} & \makecell[r]{+26.4\,\%\\{}[23.8,\,29.6]\\$p$ $<$\,0.001} & \makecell[r]{+11.9\,\%\\{}[9.6,\,14.9]\\$p$ $<$\,0.001} & \makecell[r]{+6.5\,\%\\{}[5.0,\,7.8]\\$p$ $<$\,0.001} & \makecell[c]{sí\\{}(1/3)} & FALLA \\
P3 & \makecell[l]{Red con\\{}rotonda,\\{}nominal} & \makecell[c]{3\\{}+ rot.} & 2 & 2150 & \makecell[r]{24.76\\{}$\pm$ 0.24} & \makecell[r]{+22.2\,\%\\{}[21.3,\,23.4]\\$p$ $<$\,0.001} & \makecell[r]{+13.4\,\%\\{}[12.3,\,14.7]\\$p$ $<$\,0.001} & \makecell[r]{$-$31.7\,\%\\{}[$-$33.2,\,$-$30.6]\\$p$ $<$\,0.001} & \makecell[c]{sí\\{}(2/3)} & FALLA \\
P4 & \makecell[l]{Red con\\{}rotonda,\\{}alta} & \makecell[c]{3\\{}+ rot.} & 2 & 2580 & \makecell[r]{33.36\\{}$\pm$ 0.99} & \makecell[r]{+16.8\,\%\\{}[15.0,\,18.6]\\$p$ $<$\,0.001} & \makecell[r]{+9.8\,\%\\{}[8.2,\,11.6]\\$p$ $<$\,0.001} & \makecell[r]{$-$32.5\,\%\\{}[$-$33.6,\,$-$31.6]\\$p$ $<$\,0.001} & \makecell[c]{no\\{}(2/3)} & \makecell[l]{PENDIENTE\\{}(no convergió;\\{}provisional:\\{}FALLA)} \\
P5 & \makecell[l]{Malla $3 \times 3$,\\{}nominal} & 9 & 2 & 4350 & \makecell[r]{30.16\\{}$\pm$ 0.09} & \makecell[r]{+41.9\,\%\\{}[41.5,\,42.5]\\$p$ $<$\,0.001} & \makecell[r]{+16.7\,\%\\{}[16.1,\,17.3]\\$p$ $<$\,0.001} & \makecell[r]{$-$48.8\,\%\\{}[$-$49.4,\,$-$48.2]\\$p$ $<$\,0.001} & \makecell[c]{sí\\{}(3/3)} & FALLA \\
P6 & \makecell[l]{Malla $3 \times 3$,\\{}alta} & 9 & 2 & 8700 & \makecell[r]{92.84\\{}$\pm$ 2.63} & \makecell[r]{+13.0\,\%\\{}[10.3,\,16.2]\\$p$ $<$\,0.001} & \makecell[r]{$-$15.0\,\%\\{}[$-$17.7,\,$-$11.6]\\$p$ $<$\,0.001} & \makecell[r]{$-$36.3\,\%\\{}[$-$44.8,\,$-$29.3]\\$p$ $<$\,0.001} & \makecell[c]{sí\\{}(1/3)} & FALLA \\
\bottomrule
\end{tabular}
\caption[Robustez: veredicto de Q-learning]{Robustez. Q-learning: veredicto por caso. Sem.: semáforos (rot.: más una rotonda); Carr.: carriles por sentido de la avenida; Dem.: demanda total. Retraso: retraso total (timeLoss más la espera para entrar a la red), media $\pm$ IC95\,\% sobre las semillas de evaluación 1001 a 1030; cambio (\%) con IC95\,\% bootstrap pareado y $p$ de Wilcoxon pareado. Versión evaluada: la versión guardada del entrenamiento mediano de los tres por la validación (semillas 999 y 1999). Cambio de la recompensa: el de la recompensa propia del modelo (ret\_retraso) frente a la del fijo, positivo = mejor. Conv.: si convergió el entrenamiento de la versión evaluada; entre paréntesis, cuántos de los tres entrenamientos convergen.}
\label{tab:robustez-veredicto-ql}
\end{table}
